% TOP_PROBLEM: {"id": 2561, "title": "Kourovka Notebook Problem 21.52", "category_id": 17, "category": {"id": 17, "name": "group_theory", "display_name": "Group Theory", "description": "Problems about groups, group actions, representations, and related algebraic structures.", "slug": "group-theory", "order_index": 17, "created_at": "2026-07-31T15:26:25.671Z"}, "status": "open", "problem_number": "KOU-21.52", "set_id": 10}
% FIRST_POSTED: 2026-10-01
% ENTRY_KIND: statement_only
% UnsolvedMath ID 2561; source catalogue code KOU-21.52
% !TeX program = lualatex
% Definition and sources only; this file is not an LLM solution attempt.
\documentclass[11pt,a4paper]{article}
\usepackage[margin=25mm]{geometry}
\usepackage{fontspec}
\setmainfont{lmroman10-regular.otf}[BoldFont=lmroman10-bold.otf,
 ItalicFont=lmroman10-italic.otf,BoldItalicFont=lmroman10-bolditalic.otf]
\usepackage{amsmath,amssymb,mathtools}
\usepackage{unicode-math}
\setmathfont{latinmodern-math.otf}
\AtBeginDocument{\let\setminus\smallsetminus}
\usepackage{xurl,hyperref}
\hypersetup{colorlinks=true,urlcolor=blue}
\setcounter{secnumdepth}{0}
\setlength{\parindent}{0pt}
\setlength{\parskip}{5pt}
\setlength{\emergencystretch}{3em}
\begin{document}
\section*{Kourovka Notebook Problem 21.52}\label{2561}
{\small UnsolvedMath ID 2561 · KOU-21.52 · Group Theory · Kourovka Notebook - New Problems, Issue 21}
\subsection{Definitions and mathematical statement}
Let $L$ be a finite nonabelian simple group and let $D$ be one conjugacy
class of involutions, meaning elements of order two. Form the complete
simple graph on vertex set $D$, and label its unordered edge $\{a,b\}$
by the integer $o(ab)$, the order of the product. This is well defined
on unordered pairs, since $ba=(ab)^{-1}$ for involutions and inverse
elements have the same order.

Its colour-preserving automorphism group is
\[
 A=\{\tau\in\operatorname{Sym}(D):
        o(\tau(a)\tau(b))=o(ab)\text{ for all distinct }a,b\in D\}.
\]
The colours are the actual integer orders; permuting colour labels
is not allowed. Let $\operatorname{Aut}(L)_D$ be the subgroup of group
automorphisms mapping $D$ to itself. Restriction to $D$ gives a homomorphism
$\operatorname{Aut}(L)_D\to A$.

Does every $\tau\in A$ extend to an automorphism of $L$? This is the
precise interpretation of the proposed containment of graph automorphisms
in $\operatorname{Aut}(L)$. The restriction homomorphism is injective:
$\langle D\rangle$ is a nontrivial normal subgroup of the simple group
$L$, and consequently is all of $L$.

Thus an extension, if it exists, is unique, and the target can be written
$A=\operatorname{Aut}(L)_D|_D$. The task must hold for every choice of
involution conjugacy class, not only for its union with other classes.
Preservation of pair-product orders is the only graph data supplied.
\subsection{Short English statement}
Does every permutation of an involution conjugacy class preserving all pair-product orders extend to an automorphism of the finite simple group?
\subsection{Sources}
\begin{itemize}
\item[S1] ULAM.AI and the UnsolvedMath Contributors, supplied problems.json, record 2561 (KOU-21.52), Kourovka Notebook - New Problems, Issue 21. Catalogue identity and source transcription; CC BY 4.0.
\url{https://huggingface.co/datasets/ulamai/UnsolvedMath}.
\item[S2] The Kourovka Notebook, 21st edition, new problems, Problem 21.52. Expanded formulation of the supplied catalogue statement.
\url{https://arxiv.org/pdf/1401.0300v47}.
\end{itemize}
\end{document}
